% GENERATED by docs/tools/docs.py from docs/ledger.yaml — do not edit.

\begin{center}
\renewcommand{\arraystretch}{1.25}
\small
\begin{tabular}{@{}>{\raggedright\arraybackslash}p{2.4cm}>{\raggedright\arraybackslash}p{8.2cm}>{\raggedright\arraybackslash}p{2.6cm}@{}}
\toprule
\textbf{Row} & \textbf{Statement} & \textbf{Status} \\
\midrule
L-EXP & Expansion preserves or enlarges semantics (modes \texttt{None}/Covariant); residual instances keep theirs. & \textsc{proved} \\
E-NORM & Named-type normalization is semantics-preserving (interface consumed by expansion). & \textsc{assumed} \\
E-ARR & Delayed array-literal forms keep their semantics under expansion. & \textsc{assumed} \\
E-CLASS-SYM & Class expansion either does not depend on existing variable bindings, or any bindings it reads cannot make cache reuse unsound. & \textsc{open} \\
E-GUARD & Recursion-guard identity expansions are never cached into long-lived substitutions. & \textsc{open} \\
E-PLACEHOLDER-ATOM & An un-annotated ambient-class field reads as an opaque per-field atom; the atom's denotation must over-approximate the arbitrary runtime value. & \textsc{open} \\
E-EXPKEY & The expansion-cache key (\texttt{ExpansionKey}: unqualified name text + arguments) identifies one definition: same-named, same-argument aliases from different modules must not collide. & \textsc{violated} \\
\bottomrule
\end{tabular}
\end{center}
